\section{Balanced Heath--Brown decompositions}
\label{bfk:section}

We give a fixed-weight consequence of the classical Heath--Brown
identity.  The identity itself is Heath--Brown's Lemma~1,
equation~(6), on printed page~1367 of \cite{BFKHeathBrown}.
Its application here uses the entireness of a nontrivial cubic Hecke
function to remove every lower term.  Quantitative smoothness then
restricts the remaining direct variable to a power of a logarithm.
This changes the lengths of the inverse variables; it does not prove
cancellation in their product.

\paragraph{Fixed arithmetic and analytic setup.}

Let \(K=\mathbf Q(\omega)\), where \(\omega^2+\omega+1=0\), and let
\(\psi\) be either faithful cubic ray character defining
\(K(\sqrt[3]{269})/K\), of conductor \((269)\).
Write \(L(s)=L_K(s,\psi)\) for the primitive Hecke function.
Its rational conductor is
\[
 N=|D_K|\,\mathrm N(269)=3\cdot269^2=217083.
\]
We use its standard entire continuation and functional equation
with completion
\begin{equation}
 \Lambda(s)=a_0^{-s}\Gamma(s)L(s),\qquad
 a_0=\frac{2\pi}{\sqrt N},\qquad
 L(s)=\epsilon a_0^{2s-1}
       \frac{\Gamma(1-s)}{\Gamma(s)}L^\vee(1-s),
 \quad |\epsilon|=1.
 \label{bfk:fe}
\end{equation}
Here \(L^\vee\) is the dual Hecke function.  The analytic background
needed below is finite order and polynomial growth in each fixed
vertical strip; see \cite[\S2.1]{BFKZaman}.

Delete the primes above \(3\) and \(269\), and put
\[
 P(s)=L_K^S(s,\psi)=H(s)L(s),\qquad
 H(s)=1-\psi(\lambda)3^{-s},\qquad \lambda=\omega-1.
\]
The prime above \(269\) already has no Euler factor in \(L\).
All ideals in what follows are nonzero integral ideals prime to \(S\).
Set
\[
 m(\mathfrak a)=\mu_K(\mathfrak a)\psi(\mathfrak a),\qquad
 D(s)=P(s)^{-1},\qquad
 M_U(s)=\sum_{\mathrm N\mathfrak a\le U}
          m(\mathfrak a)(\mathrm N\mathfrak a)^{-s}.
\]
Let \(d_{j,K}(\mathfrak c)\) count ordered factorizations into \(j\)
ideals.  In particular
\(d_{j,K}(\mathfrak p^v)=\binom{v+j-1}{j-1}\).
The complete multiplicativity of \(\psi\) gives
\begin{equation}
 F_j(s):=-\frac{(P^j)'(s)}{j}
   =\sum_{\mathfrak c}
     \psi(\mathfrak c)d_{j,K}(\mathfrak c)
       \frac{\log\mathrm N\mathfrak c}{j}
       (\mathrm N\mathfrak c)^{-s}
 \quad(\Re s>1).
 \label{bfk:direct}
\end{equation}
We write \(A_{j,W}(Y)\) for the corresponding sum with
the additional factor \(W(\mathrm N\mathfrak c/Y)\).

Fix \(W\in C_c^\infty((1,2))\) and let
\(\mathcal D=u\,d/du\).
Two useful assumptions on this fixed weight are
\begin{align}
 \|\mathcal D^nW\|_\infty
  &\le C_WH_W^n(n!)^\sigma,\qquad \sigma>1,
  \label{bfk:gevrey}\\
 \|\mathcal D^nW\|_\infty
  &\le C_WH_W^n n!\{\log(e+n)\}^{rn},\qquad r>1.
  \label{bfk:near}
\end{align}
Let \(0\le\chi\le1\) be a fixed smooth function on \([0,\infty)\),
equal to \(1\) on \([0,1]\) and \(0\) on \([2,\infty)\).
The cutoff \(\chi\) need not satisfy either derivative condition.
The integer \(k\ge1\) is fixed independently of \(X\).

\paragraph{Balanced prime-convolution theorem.}
Put \(U=(2X)^{1/k}\) for \(X\ge2\).  Under
\eqref{bfk:gevrey}, take
\[
 L_X=\{\log(3X)\}^{\kappa},\qquad \kappa>2k\sigma.
\]
Under \eqref{bfk:near}, take
\begin{equation}
 L_X=\{\log(3X)\}^{2k}
       \{\log(e+\log(3X))\}^{\kappa},
 \qquad \kappa>2kr.
 \label{bfk:length}
\end{equation}
Define
\begin{equation}
\begin{split}
 R_{k,W}(X;L_X)
 ={}&\frac{(-1)^{k-1}}{k}
 \sum_{\substack{\mathrm N\mathfrak c<2L_X\\
                 \mathrm N\mathfrak d_i\le U\ (1\le i\le k)}}
 \psi(\mathfrak c)d_{k,K}(\mathfrak c)\log\mathrm N\mathfrak c
       \prod_{i=1}^k m(\mathfrak d_i)\\
 &\quad{}\cdot
 W\left(\frac{\mathrm N\mathfrak c
                  \prod_i\mathrm N\mathfrak d_i}{X}\right)
 \left[1-\chi\left(\frac{L_X\prod_i\mathrm N\mathfrak d_i}{X}\right)\right].
\end{split}
\label{bfk:remainder}
\end{equation}
Then, for every \(M>0\),
\begin{equation}
 T_W(X):=
 \sum_{\mathfrak a}\Lambda_K(\mathfrak a)\psi(\mathfrak a)
                    W(\mathrm N\mathfrak a/X)
 =R_{k,W}(X;L_X)+O_{W,k,\chi,M}(X^{-M}).
 \label{bfk:main}
\end{equation}
Every nonzero summand in \eqref{bfk:remainder} satisfies
\begin{equation}
 \mathrm N\mathfrak c<2L_X,\qquad
 \frac{U}{2L_X}<\mathrm N\mathfrak d_i\le U
 \quad(1\le i\le k).
 \label{bfk:balance}
\end{equation}
These assertions are unconditional.  For any prescribed
\(\kappa>2k\) in \eqref{bfk:length}, one may choose
\(1<r<\kappa/(2k)\); nonzero nonnegative admissible weights exist
inside every open subinterval of \((1,2)\).
All constants may depend on the fixed conductor and on \(k\);
no assertion uniform in a growing \(k\) or conductor is intended.

\paragraph{Rational-prime and RH criteria.}
For \(p\ne3,269\), put
\[
 t(p)=\#\{z\bmod p:z^3=269\}-1,\qquad
 \Theta_W(X)=\sum_{\substack{p\ne3,269\\p\equiv1\ (3)}}
                t(p)\log p\,W(p/X).
\]
Then
\[
 T_W(X)-\Theta_W(X)\ll_W X^{1/2}\{\log(3X)\}^2.
\]
For either fixed weight class above, with its corresponding strict
cutoff inequality, the following are equivalent:
\begin{enumerate}
\item RH for the primitive cubic Hecke function \(L(s)\);
\item for every nonnegative weight in that class and every
\(\varepsilon>0\),
\(R_{k,W}(X;L_X)=O_{W,k,\varepsilon}(X^{1/2+\varepsilon})\);
\item for every nonnegative weight in that class and every
\(\varepsilon>0\),
\(\Theta_W(X)=O_{W,\varepsilon}(X^{1/2+\varepsilon})\).
\end{enumerate}
More generally, a bound \(\Re\rho\le\beta\) for every nontrivial
zero, with fixed \(1/2\le\beta<1\), implies
\(R_{k,W}=O_{W,k}(X^\beta)\).
The equivalence is for all the indicated localized weights, not one
selected test.  Central and repeated zeros are allowed.

\paragraph{Full inverse variant.}
Let \(k\ge2\), keep \(U=(2X)^{1/k}\), and put
\[
 \mathcal M_W(X)=\sum_{\mathfrak a}m(\mathfrak a)
                                      W(\mathrm N\mathfrak a/X).
\]
Define
\begin{equation}
\begin{split}
 \mathcal R_{k,W}(X;L_X)
 ={}&(-1)^{k-1}
 \sum_{\substack{\mathrm N\mathfrak c<2L_X\\
                 \mathrm N\mathfrak d_i\le U\ (1\le i\le k)}}
 \psi(\mathfrak c)d_{k-1,K}(\mathfrak c)
                   \prod_i m(\mathfrak d_i)\\
 &\quad{}\cdot
 W\left(\frac{\mathrm N\mathfrak c
                   \prod_i\mathrm N\mathfrak d_i}{X}\right)
 \left[1-\chi\left(\frac{L_X\prod_i\mathrm N\mathfrak d_i}{X}\right)\right].
\end{split}
\label{bfk:inverse-remainder}
\end{equation}
Then
\begin{equation}
 \mathcal M_W(X)=\mathcal R_{k,W}(X;L_X)+O_M(X^{-M})
 \label{bfk:inverse-main}
\end{equation}
for every \(M\), with the same balanced support
\eqref{bfk:balance}, if either
\[
 L_X=\{\log(3X)\}^{\kappa},\quad \kappa>2(k-1)\sigma,
\]
or
\[
 L_X=\{\log(3X)\}^{2(k-1)}
       \{\log(e+\log(3X))\}^{\kappa},
 \quad\kappa>2(k-1)r,
\]
under the respective derivative condition.



The constants may depend on the displayed fixed parameters.
These inverse identities and the balanced prime-convolution theorem
use neither cubic RH nor classical zeta RH.  No estimate for the
remaining signed sum is asserted.

\section{Proof overview}
The proof starts with Heath--Brown's binomial identity on the ideal
monoid.  The tail of the inverse series is supported beyond \(U\),
so its \(k\)-fold convolution cannot meet the shell of norm at most
\(2X=U^k\).  After multiplication by \(-P'\), the direct function
in the \(j\)-th term is exactly \(-(P^j)'/j\).  Its entireness
makes every term with \(j<k\) rapidly decreasing, because its
direct scale is a positive fixed power of \(X\).

For the last term we first prove a contour estimate whose constants
are uniform in the distance shifted to the left.  Optimizing this
estimate against the fixed derivative bounds gives direct decay of
order \(\exp\{-cY^{1/(2k)}/(\log Y)^r\}\).  The cutoff is placed
outside this direct sum, so no changing inner weight has to satisfy
new derivative bounds.  Its complement forces every inverse norm
into \((U/(2L_X),U]\).  The prime criterion is then an exact
Mellin-continuation argument; it does not supply cancellation in the
balanced expression.  Omitting \(-P'\) changes the last direct
power to \(P^{k-1}\), giving the smaller degree in the inverse
variant.

\section{Contour estimates with constants uniform in the shift}

The uniform estimate that permits optimization is the following.
Put
\[
 \mathcal B=\max\left(1,\frac{3N}{4\pi^2}\right).
\]
For every fixed \(j\ge1\), all integers \(A\ge1\), and all \(Y\ge1\),
\begin{equation}
 |A_{j,W}(Y)|
 \le C_j\mathcal B^{jA}Y^{-A}
       \|\mathcal D^{\,2jA+j+3}W\|_\infty.
 \label{bfk:uniform}
\end{equation}
The constant \(C_j\) is independent of \(A,Y,W\).
Thus the dependence on the fixed conductor in the exponential
shift constant is displayed explicitly.

To prove this, set \(G(s)=\Gamma(1-s)/\Gamma(s)\) and
\(R=(A^2+t^2)^{1/2}\).  Gamma recurrence and conjugation give
the exact equality
\begin{equation}
 |G(-A+it)|=|t|\prod_{\nu=1}^{A}(\nu^2+t^2)
       \le R^{2A+1}.
 \label{bfk:gamma}
\end{equation}
It holds also at \(t=0\), where both sides of the equality vanish.
Reflection gives
\[
 G(s)=\frac{\sin\pi s}{\pi}\Gamma(1-s)^2.
\]
For \(t\ne0\),
\[
 G'(s)=G(s)\{\pi\cot\pi s-2\operatorname{digamma}(1-s)\}.
\]
On the line in question,
\[
 |t\coth\pi t|\ll1+|t|,\qquad
 |\operatorname{digamma}(1+A-it)|\ll\log(3+R).
\]
The latter bound is uniform because its argument remains in the
right half-plane with real part at least two.  The gamma formulas
and this asymptotic are given in \cite[\S\S5.5, 5.11]{BFKDLMF}.
Multiplication before taking absolute values, followed by the
removable limit at \(t=0\), yields
\begin{equation}
 |G'(-A+it)|\ll R^{2A+1}\log(3+R).
 \label{bfk:gamma-derivative}
\end{equation}
In particular, no division by a trivial zero is used.

On \(\Re z\ge2\), absolute convergence bounds \(L^\vee(z)\) and
\((L^\vee)'(z)\) by constants depending only on \(K\).
Also
\[
 |H(-A+it)|+|H'(-A+it)|\ll 3^A.
\]
Differentiate \eqref{bfk:fe} and use
\eqref{bfk:gamma}--\eqref{bfk:gamma-derivative}.
The result is
\begin{align*}
 |P(-A+it)|&\le C\mathcal B^A R^{2A+1},\\
 |P'(-A+it)|&\le C\mathcal B^A
                    R^{2A+1}\log(3+R).
\end{align*}
The factor \(2\log a_0\) from differentiating the conductor term
is included in the fixed constant \(C\).
Since \(F_j=-P'P^{j-1}\), this proves
\begin{equation}
 |F_j(-A+it)|\le C_j\mathcal B^{jA}
                     R^{2jA+j}\log(3+R).
 \label{bfk:power-bound}
\end{equation}
There is no unrecorded factorial in this inequality.

With \(w(v)=W(e^v)\), Mellin transformation and integration by
parts against the full complex parameter \(s\) give
\begin{equation}
 |\widehat W(-A+it)|
 \le(\log2)\|\mathcal D^\ell W\|_\infty
            |-A+it|^{-\ell}\qquad(\ell\ge0).
 \label{bfk:mellin}
\end{equation}
The support lies in \(0<v<\log2\), so \(e^{-Av}\le1\), and there
are no boundary terms.  In particular, a spurious factor
\(A^\ell\) does not occur.

Mellin inversion initially on \(\Re s=2\) reads
\[
 A_{j,W}(Y)=\frac1{2\pi i}
          \int_{(2)}F_j(s)\widehat W(s)Y^s\,ds.
\]
For each fixed \(A\), entireness of \(F_j\), its polynomial bounds
in the fixed intervening strip, and rapid Mellin decay justify
moving the contour to \(\Re s=-A\), without residues.
Polynomial strip bounds for its derivative follow also by
Cauchy's formula on a slightly larger fixed strip.
The horizontal integrals tend to zero before \(A\) is allowed to
depend on \(Y\).
On the new line use \(\ell=2jA+j+3\).  The remaining integral is
bounded by a constant times
\[
 \int_{\mathbf R}(A^2+t^2)^{-3/2}
             \log(3+\sqrt{A^2+t^2})\,dt,
\]
which is uniformly bounded for \(A\ge1\) after substituting
\(t=Av\).  This proves \eqref{bfk:uniform}.
Only the last vertical estimate, already uniform in \(A\), is
used when optimizing the shift.

\paragraph{Optimizing the derivative classes.}
Under \eqref{bfk:gevrey}, \eqref{bfk:uniform} and \(n!\le n^n\)
give constants \(C_1,C_2,C_3\), independent of \(A,Y\), such that
\[
 |A_{j,W}(Y)|
 \le C_1(A+1)^{C_2}
       \left\{\frac{C_3(A+1)^{2j\sigma}}Y\right\}^{A}.
\]
Choose \(A=\lfloor\delta Y^{1/(2j\sigma)}\rfloor\), with fixed
sufficiently small \(\delta>0\).  The expression in braces is
bounded away from one for large \(Y\), and the polynomial factor
is absorbed by the exponential.  Enlarging the constant on a
bounded initial range gives
\begin{equation}
 |A_{j,W}(Y)|\ll_{j,W}
        \exp(-c_{j,W}Y^{1/(2j\sigma)}).
 \label{bfk:gevrey-decay}
\end{equation}
For \eqref{bfk:near}, the corresponding bound has the form
\[
 C_1(A+1)^{C_2}\{\log(e+A)\}^{C_3}
 \left\{\frac{C_4(A+1)^{2j}
                    \{\log(e+A)\}^{2jr}}Y\right\}^{A}.
\]
Taking
\[
 A=\left\lfloor
       \delta\frac{Y^{1/(2j)}}{\{\log(3Y)\}^{r}}
    \right\rfloor
\]
instead yields
\begin{equation}
 |A_{j,W}(Y)|\ll_{j,W}
 \exp\left(-c_{j,W}
             \frac{Y^{1/(2j)}}{\{\log(3Y)\}^{r}}\right).
 \label{bfk:near-decay}
\end{equation}
The power of the logarithm in the denominator remains \(r\);
the factor \(j\) enters the root \(1/(2j)\).
The constants in these estimates are uniform over a family with
fixed bounds \(C_W,H_W,\sigma\), or \(C_W,H_W,r\).
They are not uniform when the weight constants or conductor vary
without restriction.

\paragraph{Localized nonnegative weights.}
Here is an explicit construction, included for the quantifier
over weights used below.  Given an open interval \(J\subset(1,2)\),
choose positive widths \(a_n\) with
\(\sum_n a_n<|\log J|\).  Convolve the uniform probability
measures on \([-a_n/2,a_n/2]\).  Their sum exists because the
widths are summable, and its distribution has support of width
at most \(\sum_n a_n\).
Differentiating the first \(m\) box measures and using the
\((m+1)\)-st for the supremum norm gives
\[
 \|w^{(m)}\|_\infty
 \le \frac{2^m}{a_1\cdots a_{m+1}}.
\]
Keeping one further box shows that every distributional
derivative is continuous: the convolution of two box densities
is continuous, and convolution with a finite measure preserves
continuity here.  The resulting density is \(C^\infty\),
nonnegative, and has integral one.
Translate it strictly inside \(\log J\), and set \(W(u)=w(\log u)\).
Widths \(a_n=c n^{-\sigma}\) give \eqref{bfk:gevrey}.
Widths
\[
 a_n=\frac{c}{n\{\log(e+n)\}^{r}},\qquad r>1,
\]
give \eqref{bfk:near}.  In both cases choose \(c>0\) small enough
to fit the interval; the extra factor at index \(m+1\) is absorbed
by a fixed exponential in \(m\).

\section{The exact identity and the balanced truncation}

In the absolutely convergent half-plane, the binomial theorem
gives
\begin{equation}
 D=\sum_{j=1}^{k}(-1)^{j-1}\binom{k}{j}P^{j-1}M_U^j
          +P^{k-1}(D-M_U)^k.
 \label{bfk:identity}
\end{equation}
Indeed \(P(D-M_U)=1-PM_U\), and
\(D[1-(1-PM_U)^k]\) is the displayed finite sum.
This proves the identity as a formal ideal convolution as well.
Each factor \(D-M_U\) has all its ideal coefficients supported
at norms strictly greater than \(U\).  Thus the last term in
\eqref{bfk:identity}, and also its product with \(-P'\), has
no coefficient at norm at most \(U^k=2X\).
Multiplication by any of the other ideal series cannot decrease
the norm.

Multiplying by \(-P'\), using
\(-P'P^{j-1}=-(P^j)'/j=F_j\), and extracting the \(W\)-weighted
shell proves the exact formula
\begin{equation}
 T_W(X)=\sum_{j=1}^k(-1)^{j-1}\binom{k}{j}
  \sum_{\mathrm N\mathfrak d_1,\ldots,\mathrm N\mathfrak d_j\le U}
   \left(\prod_{i=1}^j m(\mathfrak d_i)\right)
   A_{j,W}\left(\frac X{\prod_i\mathrm N\mathfrak d_i}\right).
 \label{bfk:exact}
\end{equation}
No differentiation of \(M_U\) is involved: \(-P'\) multiplies
the identity.  This is the source of the exact coefficient
\((\log\mathrm N\mathfrak c)/j\).
All character values and signs have been kept.

The number of ideals of \(K\) of norm at most \(Z\) is \(O(Z)\).
For an elementary proof, the PID property supplies six generators
per ideal, and
\[
 a^2-ab+b^2\ge\tfrac12(a^2+b^2)
\]
bounds the relevant lattice points by \(O(Z)\).
Consequently there are \(O_j(U^j)\) tuples in the \(j\)-th term.
For \(j<k\), their direct scales satisfy
\[
 \frac X{\prod_i\mathrm N\mathfrak d_i}
 \ge 2^{-j/k}X^{1-j/k}.
\]
For sufficiently large \(X\), these scales are at least one.
Choose a fixed integer \(A\) in \eqref{bfk:uniform}, sufficiently
large in terms of \(M,j,k\).  Summing that bound gives
\[
 O_{W,j,k,A}\left(X^{\,j/k-A(1-j/k)}\right)
   =O_{W,k,M}(X^{-M}).
\]
The bounded initial range of \(X\) is absorbed in the constant.
Thus every lower term is negligible to arbitrary order even for
an arbitrary fixed \(C^\infty\) weight.  When \(k=1\) there are no
lower terms.  The last term of \eqref{bfk:exact} has coefficient
\((-1)^{k-1}\), since \(\binom{k}{k}=1\).

Put \(b=\prod_i\mathrm N\mathfrak d_i\) in this last term, and
split it using \(\chi(L_X b/X)\).
The omitted part is bounded by
\[
 \sum_{\mathrm N\mathfrak d_i\le U}
    \left|\chi(L_X b/X)A_{k,W}(X/b)\right|.
\]
There are \(O_k(U^k)=O_k(X)\) tuples.  Where this cutoff is
nonzero, \(X/b\ge L_X/2\).  Therefore
\eqref{bfk:gevrey-decay} makes the omitted part
\begin{equation}
 O_{k,W}\left(
       X\exp\{-c(L_X/2)^{1/(2k\sigma)}\}\right).
 \label{bfk:gevrey-error}
\end{equation}
For the nearly analytic class, monotonicity of
\(Y^{1/(2k)}/\{\log(3Y)\}^r\) for large \(Y\) similarly gives
\begin{equation}
 O_{k,W}\left(
       X\exp\left\{-c\frac{L_X^{1/(2k)}}
                              {\{\log(3L_X)\}^r}\right\}\right).
 \label{bfk:near-error}
\end{equation}
Only the bound \(|\chi|\le1\) was used.  In particular no
derivative of a changing cutoff or an inner test function enters
these estimates.

For the Gevrey length, the exponent in
\eqref{bfk:gevrey-error} divided by \(\log X\) tends to infinity
precisely under the strict inequality used in the theorem.
For \eqref{bfk:length}, put \(t=\log(3X)\) and \(h=\log(e+t)\).
Then \(\log(3L_X)\ll_{k,\kappa}h\), and the positive exponent
in \eqref{bfk:near-error} is bounded below by
\[
 c_{k,\kappa}t\,h^{\kappa/(2k)-r}.
\]
Its ratio to \(\log X\) tends to infinity when \(\kappa>2kr\).
Thus both omitted parts are \(O_M(X^{-M})\) for every \(M\).
The bounded initial range of \(X\) is absorbed in the constant.
At either equality endpoint this reasoning supplies at most a
fixed power saving determined by the weight's decay constant,
not arbitrary powers.

In the retained part, \(b>X/L_X\), whereas the support of \(W\)
implies \(\mathrm N\mathfrak c\,b<2X\).  Hence
\(\mathrm N\mathfrak c<2L_X\).
For each \(i\), the other \(k-1\) factors are at most \(U^{k-1}\),
and
\[
 \mathrm N\mathfrak d_i>
       \frac{X}{L_XU^{k-1}}=\frac{U}{2L_X}.
\]
This proves \eqref{bfk:remainder}--\eqref{bfk:balance}.
It also verifies the \(k=1\) specialization: the explicit upper
bound \(2X\) on its sole inverse variable is already forced by
the shell support, and the usual switched convolution results.

\section{Proof of the rational-prime criterion}

For \(p\ne3,269\), let
\[
 t(p)=\#\{z\bmod p:z^3=269\}-1,\qquad
 \Theta_W(X)=
 \sum_{\substack{p\ne3,269\\p\equiv1\ (3)}}
       t(p)\log p\,W(p/X).
\]
At a split prime the two character values over \(p\) sum to
\(t(p)\), by the cubic Kummer character.  At an inert prime the
residue degree is two, so it gives no degree-one term.
Consequently \(\Theta_W\) is exactly the degree-one,
first-power part of \(T_W\).
All remaining terms have rational norm \(p^v\), \(v\ge2\).
There are at most two prime ideals over each rational prime,
and their logarithmic weights are at most \(2\log p\).
Bounding the number of possible \(p\)'s by
\((2X)^{1/v}\), and the number of exponents by \(O(\log(3X))\),
proves the elementary estimate
\begin{equation}
 T_W(X)-\Theta_W(X)
      \ll_W X^{1/2}\{\log(3X)\}^2.
 \label{bfk:prime-powers}
\end{equation}
No prime number theorem or zero hypothesis is needed for this
comparison.

\paragraph{Equivalent square-root criteria.}
For either fixed class \eqref{bfk:gevrey} or \eqref{bfk:near},
with the corresponding strict cutoff inequality, the following
are equivalent:
\begin{enumerate}
\item RH for the primitive cubic Hecke function \(L(s)\);
\item for every nonnegative weight in that class and every
      \(\varepsilon>0\),
      \(R_{k,W}(X;L_X)=O_{W,k,\varepsilon}(X^{1/2+\varepsilon})\);
\item for every such weight and every \(\varepsilon>0\),
      \(\Theta_W(X)=O_{W,\varepsilon}(X^{1/2+\varepsilon})\).
\end{enumerate}
More generally, if every nontrivial zero has real part at most
\(\beta\), where \(1/2\le\beta<1\), then
\begin{equation}
 R_{k,W}(X;L_X)=O_{W,k}(X^\beta).
 \label{bfk:zero-bound}
\end{equation}

We include the zero argument to specify its scope.
The Hadamard logarithmic derivative and gamma bounds give
\[
 \#\{\rho:|\Im\rho-t|\le1\}\ll\log(3+|t|),
\]
with multiplicities included; see
\cite[Lemmas~2.2--2.3]{BFKZaman}.
Indeed, at \(2+it\) the real zero summands are nonnegative,
and each nearby zero contributes at least a fixed positive
constant, while the completed logarithmic derivative is
\(O(\log(3+|t|))\).
Subtracting that Hadamard identity from the one at \(\sigma+it\)
shows, in a fixed strip, that the contribution of zeros farther
than one in ordinate is \(O(\log(3+|t|))\): the subtracted
reciprocal terms have a summable inverse-square tail.
Choose heights tending to infinity at distance
\(\gg1/\log(3+|t|)\) from all neighboring zero ordinates.
The local count ensures such heights exist.  The logarithmic
derivative on those horizontal contours is
\(O(\log^2(3+|t|))\).

Shift the Mellin integral for the primitive logarithmic
derivative from \(2\) to \(-1/2\).
Rapid Mellin decay removes the horizontal integrals.
On the new line the functional equation and the absolutely
convergent dual logarithmic derivative at real part \(3/2\)
give logarithmic vertical growth.  Nontrivial zeros contribute
their full residues \(-m_\rho\); the possible residue at zero
is bounded independently of \(X\).
Removing the one extra Euler factor changes the prime-power
sum by \(O_W(1)\), because only a bounded number of powers of a
fixed prime can lie in \((X,2X)\).
It follows that
\begin{equation}
 T_W(X)=-\sum_\rho m_\rho\widehat W(\rho)X^\rho+O_W(1).
 \label{bfk:zeros}
\end{equation}
The sum is absolutely convergent: the local count and rapid
decay of \(\widehat W(\sigma+it)\), uniformly for \(0\le\sigma\le1\),
give \(\sum_\rho m_\rho|\widehat W(\rho)|<\infty\).
Thus the asserted zero bound proves \eqref{bfk:zero-bound} through
\eqref{bfk:main}; under RH it even gives
\(R_{k,W}=O_{W,k}(X^{1/2})\).
A central zero of multiplicity \(m\) contributes exactly
\(-m\widehat W(1/2)X^{1/2}\).  No central zero, repeated zero,
or off-line zero was discarded.

For the converse, either proposed square-root criterion,
together with \eqref{bfk:main} and \eqref{bfk:prime-powers},
gives the same \(X^{1/2+\varepsilon}\) bound for \(T_W\).
The Mellin integral
\[
 \int_0^\infty T_W(X)X^{-s-1}\,dX
\]
is then holomorphic for \(\Re s>1/2\), locally uniformly.
It vanishes near \(X=0\), and in \(\Re s>1\) it equals
\begin{equation}
 -\widehat W(s)\frac{P'}{P}(s).
 \label{bfk:criterion}
\end{equation}
If \(L\) had a zero \(\rho\) to the right of \(1/2\), choose
a sufficiently short open interval \(J\subset(1,2)\) so that
the phases \(u^{i\Im\rho}\) on \(J\) lie in an open half-plane
after a fixed rotation.
The constructed class contains a nonzero nonnegative weight
supported in \(J\), and hence \(\widehat W(\rho)\ne0\).
The Euler polynomial \(H\) has no zero in \(\Re s>0\).
Expression \eqref{bfk:criterion} would therefore have a pole
at \(\rho\) of nonzero residue, a contradiction.
Functional-equation symmetry, together with conjugation of
the character, excludes the reflected zeros to the left.
This proves RH and the equivalences.
The quantifier over all localized weights is essential:
one particular Mellin transform may vanish at a specified zero.

\section{Proof of the undifferentiated inverse variant}

For completeness, an analogous result with a smaller direct
degree follows from \eqref{bfk:identity} itself.
Let \(k\ge2\) and
\[
 \mathcal M_W(X)=\sum_{\mathfrak a}
             m(\mathfrak a)W(\mathrm N\mathfrak a/X).
\]
For each \(j\ge2\) the direct function is \(P^{j-1}\), not
the derivative of \(P^j\).
The \(j=1\) term is \(kM_U\) and has empty shell support,
since \(U\le X\) for \(X\ge2\).
The residual has empty support as before.
The same contour proof, without the logarithmic derivative,
gives for the direct coefficients of \(P^q\)
\begin{equation}
 \left|\sum_{\mathfrak c}\psi(\mathfrak c)d_{q,K}(\mathfrak c)
                    W(\mathrm N\mathfrak c/Y)\right|
 \le C_q\mathcal B^{qA}Y^{-A}
               \|\mathcal D^{\,2qA+q+2}W\|_\infty
 \quad(q\ge1).
 \label{bfk:inverse-direct}
\end{equation}
Here the remaining vertical integral is bounded by
\(\int_{\mathbf R}(A^2+t^2)^{-1}\,dt=\pi/A\).
Thus all terms \(j<k\) are rapidly decreasing, and the last
term has direct degree \(2(k-1)\).

Equations \eqref{bfk:inverse-remainder}--\eqref{bfk:inverse-main}
now follow by the same cutoff argument, using
\eqref{bfk:inverse-direct} with \(q=k-1\).

\paragraph{What the decomposition does not estimate.}
For \(k=2\), \eqref{bfk:inverse-main} leaves two inverse variables
of size \(X^{1/2}\) up to logarithms and a direct variable of
length \((\log X)^2(\log\log X)^\kappa\), \(\kappa>2\),
after choosing an admissible fixed \(r\).
The prime-convolution version at \(k=2\) instead has direct
length \((\log X)^4(\log\log X)^\kappa\), \(\kappa>4\).
For general fixed \(k\), the trade-off is analogous.
No factor in the remaining product is replaced by its absolute
value in the identity, and no estimate for its signed sum follows
from the support restrictions.
These are applications of the classical identity with explicit
regular-weight estimates, not a new combinatorial identity or
a new zero-free result.
Finally, the primitive-content series is \(C(s)/P(s)\); its
reciprocal \(P(s)/C(s)\) is not supplied as an entire direct Hecke
function.  The undifferentiated argument here concerns the full
ideal inverse series and does not silently apply to that
primitive-content series.
